\chapter[Siegel-disk renormalization]{Source digest: golden-mean Siegel disks and renormalization}
\label{ch:digests-siegel}

This chapter digests the paper of Gaidashev and Yampolsky on golden-mean Siegel disks \cite{GaYaGoldenMean}.  It belongs to
the Siegel-disk neighbourhood of the approaches to MLC (Chapter~\ref{ch:approaches-map}), but it is not a result about the
local connectivity of $\M$: it concerns the boundaries of Siegel disks of quadratic polynomials with golden-mean rotation
number, and the renormalization that describes them.  The conventions are those of the earlier digest chapters.

\section{Setting}
\label{sec:siegel-setting}

\paragraph{Identity.}
D.~Gaidashev and M.~Yampolsky, \emph{Golden mean Siegel disk universality and renormalization}, arXiv:1604.00717 (2016).  The
title and authors are confirmed by the metadata and by \texttt{refs/arxiv-1604.00717/main.tex}, lines 203--209.  The statements
are in the files \texttt{intro.tex}, \texttt{symmetric.tex}, \texttt{circle.tex}, \texttt{dissipative.tex}, which the main file
inputs.

\paragraph{Definitions.}
\begin{itemize}
\item \emph{The golden-mean polynomial.}  $P_\theta(z)=z^2+e^{2\pi i\theta}z$ with $\theta=(\sqrt5-1)/2$ (\texttt{intro.tex},
lines 4--7).
\item \emph{The Siegel disk.}  $\Delta_\theta$ is the maximal neighbourhood of $0$ in which a conformal change of coordinates
linearizes $P_\theta$ to $w\mapsto e^{2\pi i\theta}w$ (\texttt{intro.tex}, lines 9--12).  By the results of Douady, Ghys, Herman and
Shishikura, its boundary is a Jordan curve through the unique critical point, so it coincides with the postcritical set
(lines 12--14).
\item \emph{Renormalization.}  Widom's renormalization blows up a part of $\partial\Delta_\theta$ near the critical point; the
renormalization fixed point is conjectured to be hyperbolic with one-dimensional unstable direction (\texttt{intro.tex}, lines
22--27).  The present paper works with almost commuting holomorphic pairs (\texttt{intro.tex}, lines 53--55).
\end{itemize}

\paragraph{History stated in the source.}
MacKay and Percival extended the conjecture to other rotation numbers (1986).  Stirnemann gave a computer-assisted proof of the
existence of a fixed point with golden-mean Siegel disk (1994).  McMullen proved asymptotic self-similarity of golden-mean Siegel
disks in the quadratic family, using holomorphic commuting pairs of de~Faria, with geometric convergence of renormalizations to a
fixed point (1998) (\texttt{intro.tex}, lines 29--41).  The cylinder renormalization of Yampolsky and the results of Inou--Shishikura
treat rotation numbers of high type, and do not cover the golden-mean case (lines 47--52).

\section{Results of the paper}
\label{sec:siegel-results}

\begin{theorem}[Computer-assisted renormalization fixed point]
\label{thm:gy-main1}
There exists a pair of polynomials $(\phi_0,\psi_0)$ in a space $\mathcal A_1(U,V)$ such that: the renormalization operator is
anti-analytic in a ball $B_r((\phi_0,\psi_0))$, with $r=6.6610992\cdot10^{-11}$; there is a pair $(\phi_*,\psi_*)$ in that ball that is
fixed by the operator; and $\phi_*$ and $\psi_*$ extend analytically and univalently to explicit domains containing $U$ and $V$
(\texttt{symmetric.tex}, lines 324--333).  \emph{Status:} proved with computer assistance; the constants are the source's numerical
bounds, and the proof is in the computer-assisted section (\texttt{intro.tex}, line 93).
\end{theorem}

\begin{theorem}[Quasi-arc postcritical sets]
\label{thm:gy-main4}
The postcritical set of an almost commuting pair in the stable manifold of the renormalization fixed point $\zeta_*$ is an
invariant quasi-arc (\texttt{circle.tex}, around line 6; the abstract).  \emph{Status:} proved in the paper.  Together with
McMullen's theory of holomorphic pairs, this gives the statement that for every map in the stable manifold of the golden-mean fixed
point the Siegel disk is a quasidisk whose boundary coincides with the closure of the critical orbit (abstract; \texttt{intro.tex}, lines 74--76).
\end{theorem}

\begin{theorem}[Two-dimensional dissipative maps]
\label{thm:gy-main6}
The point $\iota(\zeta_*)$ is a fixed point of $\mathcal R^n$, where $n$ is an even number and $\mathcal R$ is the renormalization of
the dissipative extension; the linear operator $N=D\mathcal R^n|_{\iota(\zeta_*)}$ is compact, and its spectrum coincides with the
spectrum of the operator $M^n$ of the paper's Main Theorem~3 (\texttt{dissipative.tex}, lines 334--339).  \emph{Status:} proved in the
paper.  The source says that these results give the hyperbolicity of renormalization for H\'enon-like maps, which is used in a
forthcoming joint paper with Radu (\texttt{intro.tex}, lines 77--80).
\end{theorem}

\section{Status and scope}
\label{sec:siegel-scope}

\paragraph{Status.}
The fixed-point and hyperbolicity conclusions for analytic one-dimensional maps rest on the computer-assisted proof of
Theorem~\ref{thm:gy-main1}.  The two-dimensional statement, Theorem~\ref{thm:gy-main6}, is proved in the paper with the same
renormalization operator.

\paragraph{MLC relevance.}
The paper contains no statement about the local connectivity of $\M$ or about the connectedness of Julia sets of a family.  Its
quasi-arc conclusion concerns the Siegel disk of individual maps in a stable manifold, and it is not a statement that $\M$ is
locally connected at any parameter.  The chapter therefore records the result in the Siegel row of the approaches map only as an
input to the Siegel-disk programme.

\section{External inputs}
\label{sec:siegel-inputs}

Siegel's linearization theorem; Douady, Ghys, Herman and Shishikura on the boundary of Siegel disks; Widom's renormalization
(1983) and the MacKay--Percival conjecture (1986); Stirnemann's computer-assisted fixed point (1994); McMullen's
self-similarity and holomorphic commuting pairs (1998); de~Faria's commuting pairs; the cylinder renormalization of Yampolsky
(critical circle maps); the Inou--Shishikura results for high-type rotation numbers; the earlier numerical paper of
Gaidashev--Yampolsky; de~Carvalho--Lyubich--Martens for the Feigenbaum renormalization on H\'enon-like maps, which the paper
extends to Siegel renormalization; and the forthcoming Gaidashev--Radu--Yampolsky paper on the Siegel disk of H\'enon-like maps
(\texttt{intro.tex}, lines 9--80).
